\section{Optical transfer between the two true ends of a short segment}
\label{sec:v84-pairs}
The preceding short-residence source can be resolved further. On a
short segment whose outgoing gradient is not of order $\sqrt r$,
all of the optical transfer between its two ends is estimated in
essential height. In particular, a grazing incoming end creates no
unestimated outgoing surface term.

For this argument we use two additional, precisely stated inherited
volume bounds. With $z=|g|$ and $I_m$ as before, let
\begin{equation}\label{eq:v84-auxiliary-sources}
 v^c(t)\,dt=F_\#(I_m z^{-1}\one_{\{z\ge c\}}\one_\Pi\nu_R^*),
 \qquad
 u^b(t)\,dt=F_\#(I_m\one_{\{z<b\}}\one_\Pi\nu_R^*).
\end{equation}
For every previously fixed $c,b,L>0$, the source-pinned estimates are
\begin{equation}\label{eq:v84-truncated-input}
 \limsup_m\sup_{R,n,k}\sup_{|J|=L}m^2\int_Jv^c\le CL,
 \qquad
 \limsup_m\sup_{R,n,k}\sup_{|J|=L}m^2\int_Ju^b\le Cb^\alpha L.
\end{equation}
The leading constant in the first bound is independent of $c>0$;
this is Supplement equation~\eqref{supp-v81archive:eq:v81-full-truncation}.
The second is Supplement Theorem~\ref{supp-v81archive:thm:v81-critical-volume},
with $\alpha=129/128$ and $b<b_1$. The first source is actually
truncated, at a positive number fixed before the collision limit.
No assertion at $c=0$ is used.

On each local angular sheet of a maximal regular segment write its
roof as $u\in(a,b)$ and define its physical angular density by
\begin{equation}\label{eq:v84-angular-density}
 A(u,\vartheta)=\frac{\omega(u,\vartheta)}
                         {4\pi c_*\chi(u,\vartheta)}.
\end{equation}
Then $d\mu_u=A(u,\vartheta)\,d\vartheta$ on that sheet. At a
nongrazing outgoing end $d\zeta_+=A(b-,\vartheta)\,d\vartheta$.
Angular sheets are disjointized only to sum their kernels; no sheet
edge is a new stopping boundary.

\begin{lemma}[Variation of the physical angular density]
\label{lem:v84-angular-variation}
Along every regular segment,
\begin{equation}\label{eq:v84-angular-variation}
 |\partial_u A(u,\vartheta)|
                    \le C A(u,\vartheta)I_m(u,\vartheta)/z(u,\vartheta).
\end{equation}
The constant is independent of collision count and all positive
internal incidences. If the incoming endpoint has positive limiting
$z$ and a vanishing internal incidence, then $A(a+,\vartheta)=0$.
If all limiting incidences and $z$ are positive, the incoming trace
exists on the corresponding one-sided physical face.
\end{lemma}
\begin{proof}
The optical sensitivity gives $|X\omega|\le C\omega I_m/z$.
Since the momentum angle is preserved,
$\chi=e_\perp^{\mathsf T}Ke_\perp$ has derivative bounded by
$T_*|X_q|\le C/z$, using the uniform cubic estimate and
$|X_q|\le C/z$. Hence $|X\log\chi|\le C/z$. Differentiate
\eqref{eq:v84-angular-density} and use $I_m\ge1$.

At fixed count, vanishing of one internal incidence forces
$\omega\to0$ by the retained selected-incidence optical identity.
The lower bound for $\chi$ proves the zero trace. Otherwise the
stationary word extends smoothly to the limiting point by the
regular reflection equations. This is a one-sided algebraic extension,
not a continuation of the physical trajectory past a first-hit or
section-decision boundary. Its positive noncritical trace is the
incoming physical flux in these angular coordinates.
\end{proof}

Partition $\mathcal S_r$ into $\mathcal S_r^{\rm nc}$ and
$\mathcal K_r$ by, respectively,
\begin{equation}\label{eq:v84-short-gradient-split}
 z(b-)^2\ge8\Lambda_\circ r,
 \qquad z(b-)^2<8\Lambda_\circ r.
\end{equation}
The equality case belongs to the first part. Let
$c_r=2\sqrt{\Lambda_\circ r}$. Since $b-a\le2r$ and
$\partial_u z^2\le2\Lambda_\circ$, every segment in the first
part has $z(u)\ge c_r$ throughout. In particular it has no critical
incoming end.

For each such segment keep the same exit-window kernel $k_r(t;y)$,
but replace its outgoing angular density $A(b-,\vartheta)$ by the
incoming angular density $A(a+,\vartheta)$. Sum this nonnegative
quantity over all angular sheets to define $\mathcal J_r(t)$.
An incoming grazing end has zero trace and contributes zero. All
remaining terms are attached to a pair of actual nongrazing incoming
and outgoing faces of the \emph{same maximal exact-event segment}.
The kernel still uses the outgoing roof $b$, not the incoming roof
$a$. Thus $\mathcal J_r$ is a paired boundary assignment, not the
unrestricted flux of all incoming faces and not a volume surrogate.

\begin{theorem}[A complete essential-height bound for boundary transfer]
\label{thm:v84-paired-transfer}
For almost every $t$,
\begin{equation}\label{eq:v84-pair-before-limit}
 |\mathcal S_r^{\rm nc}(t)-\mathcal J_r(t)|
                      \le C\int_{t-2r}^{t+2r}v^{c_r}(u)\,du.
\end{equation}
Both assignments are locally finite almost everywhere at fixed count.
Consequently
\begin{equation}\label{eq:v84-pair-height}
 \limsup_m\sup_{R,n,k}m^2
       \|\mathcal S_r^{\rm nc}-\mathcal J_r\|_\infty\le Cr.
\end{equation}
The complete part of $\mathcal S_r^{\rm nc}$ whose incoming end
is grazing satisfies the same bound $Cr$ by itself.
\end{theorem}
\begin{proof}
On one selected segment the variation inequality gives
\[
 |A(b-,\vartheta)-A(a+,\vartheta)|
       \le C\int_a^b A(u,\vartheta)I_m(u,\vartheta)/z(u,\vartheta)\,du.
\]
First use compact subsegments and then take the one-sided limits.
Multiply by $k_r(t;y)\le1$. Whenever it is nonzero, $a<t<b$;
since $b-a\le2r$, the full segment lies in $(t-2r,t+2r)$.
At each roof $u$, the map from selected outgoing faces to their
points on the same segment is injective in the physical collision
source. Indeed two such faces with the same interior state would
be the same first exit, by uniqueness of the flow. Different visits
to the same angle on different physical sheets remain different
states and are already included in $\mu_u$.

The complete sum of the variation integrals is therefore bounded by
the source in \eqref{eq:v84-auxiliary-sources}, with $z\ge c_r$.
Tonelli and the common physical kernels prove
\eqref{eq:v84-pair-before-limit}, without a word-count factor.
The source $v^{c_r}$ is finite at every fixed count. The outgoing
assignment is locally integrable by Theorem~\ref{thm:v83-surface-local-mass}
and its fixed-count proof. The variation bound thus also proves local
integrability of the incoming paired assignment; taking a difference
has not subtracted two undefined infinities. Apply
\eqref{eq:v84-truncated-input} with interval length $4r$ to obtain
\eqref{eq:v84-pair-height}. Here $c_r$ depends only on the previously
fixed outer width, not on collision count. Finally restrict the same
positive variation argument to incoming grazing ends, where the
incoming trace is zero.
\end{proof}

\begin{proposition}[Quantitative mass of the short critical part]
\label{prop:v84-critical-mass}
Put $b_r=\sqrt{8\Lambda_\circ r}<b_1$. For fixed $r,L>0$,
\begin{align}
 \limsup_m\sup_{R,n,k}\sup_{|J|=L}m^2\int_J\mathcal K_r(t)\,dt
 \le{}& C Lb_r^\alpha\notag\\
       &+C_\eta(L+r)r^{(1+\eta)/2}b_r^{\alpha-\eta}.
                                           \label{eq:v84-critical-mass}
\end{align}
In particular the mass on intervals of length $r$ is
$O_\eta(r^{385/256})$. This statement does not bound the essential
height of $\mathcal K_r$.
\end{proposition}
\begin{proof}
The radial monotonicity implies $z(u)<b_r$ on every selected segment.
Use the positive-variation direction of the proof of
Theorem~\ref{thm:v83-two-sided-flux} on just these first exits.
Its starting term is bounded by the unweighted exact source with
$z<b_r$, which is at most $u^{b_r}$ since $I_m\ge1$.
Its drift term is bounded by
$C_\eta\mathsf P_{\eta,r}v^{\eta,b_r}$: all intermediate targets
of the selected segment still have $z<b_r$. The full source on the
upper side is only enlarged, not redefined. Integrate over $J$,
use Fubini and the bounds \eqref{eq:v84-truncated-input} and
\eqref{eq:v83-fixed-width-input}. This proves the display. For
$L=r$, its exponents are $1+\alpha/2=385/256$ and
$1+(1+\alpha)/2=513/256$, respectively.
\end{proof}

Write $\mathfrak J_M(r)$ and $\mathfrak K_M(r)$ for the ordered
central essential heights of $\mathcal J_r$ and $\mathcal K_r$.
The preceding theorems imply
\begin{equation}\label{eq:v84-surface-reduction}
 \mathfrak F_M^{\rm all}(r)
       \le\mathfrak J_M(r)+\mathfrak K_M(r)+C_\eta r^{1/32}.
\end{equation}
Conversely $\mathfrak K_M(r)\le\mathfrak F_M^{\rm all}(r)$ and
$\mathfrak J_M(r)\le\mathfrak F_M^{\rm all}(r)+Cr$.
Thus full surface-height decay is equivalent to decay of the two
specified short-segment heights. Neither is inferred from its local
mass. The newly controlled terms are all long-residence outgoing
faces, all noncritical short grazing entrances, and the complete
optical transfer between the two noncritical ends of every short
segment.
